% ============================================================================
\chapter{The Tsetlin Automaton}
\label{ch:automaton}
% ============================================================================

The Tsetlin Machine is built from a single, remarkably simple
component: the \concept{Tsetlin Automaton}. This chapter introduces the
automaton, explains its behaviour, and shows why it can learn from
noisy feedback without any calculus.


% ----------------------------------------------------------------------------
\section{Finite-State Learning Automata}
\label{sec:automata-intro}
% ----------------------------------------------------------------------------

A \concept{learning automaton} is a device that repeatedly chooses an
action, receives feedback from its environment, and updates its
internal state to improve future choices. The idea dates to the 1960s
and the work of M.~L.~Tsetlin~\cite{tsetlin1961}.

The simplest version has:
\begin{itemize}
  \item A finite set of \emph{states} (numbered 1 through $2N$ for
    some positive integer $N$).
  \item Two \emph{actions}: action~0 and action~1.
  \item A rule that maps the current state to an action.
  \item An update rule that moves the state up or down based on
    feedback (reward or penalty).
\end{itemize}

\begin{analogy}
Think of a physical slider on a rail, marked with positions 1 through
$2N$. The slider starts somewhere in the middle. A ``reward'' pushes
it toward the nearest end; a ``penalty'' pushes it toward the centre.
Once the slider reaches an end, it tends to stay there. The two ends
represent two different decisions, and the slider's position reflects
the automaton's confidence in that decision.
\end{analogy}


% ----------------------------------------------------------------------------
\section{The Two-Action Tsetlin Automaton}
\label{sec:two-action}
% ----------------------------------------------------------------------------

\begin{definition}[Two-action Tsetlin Automaton with $2N$ states]
\label{def:tsetlin-automaton}
The automaton has states $\{1, 2, \ldots, 2N\}$ and two actions:
\begin{equation}
  \text{action}(s) = \begin{cases}
    0 & \text{if } s \leq N \\
    1 & \text{if } s > N
  \end{cases}
  \label{eq:action}
\end{equation}

States 1 through $N$ map to action~0 (the ``exclude'' action).
States $N{+}1$ through $2N$ map to action~1 (the ``include'' action).
\end{definition}

The state transition rules depend on the feedback signal:

\begin{description}[leftmargin=1em]
  \item[Reward:] Move the state away from the decision boundary
    (toward the nearest end of the rail):
    \begin{equation}
      s \leftarrow \begin{cases}
        \max(1, \; s - 1) & \text{if } s \leq N \\
        \min(2N, \; s + 1) & \text{if } s > N
      \end{cases}
      \label{eq:reward}
    \end{equation}

  \item[Penalty:] Move the state toward the decision boundary
    (toward the centre):
    \begin{equation}
      s \leftarrow \begin{cases}
        \min(N, \; s + 1) & \text{if } s \leq N \\
        \max(N{+}1, \; s - 1) & \text{if } s > N
      \end{cases}
      \label{eq:penalty}
    \end{equation}
\end{description}

Note that the penalty can push the state across the boundary (from
$N{+}1$ to $N$ or vice versa), changing the action. This is how the
automaton changes its mind.

\begin{figure}[htbp]
\centering
\begin{tikzpicture}[
  state/.style={circle, draw, minimum size=0.9cm, font=\small\ttfamily,
    inner sep=0pt, line width=0.6pt},
  action0/.style={state, fill=bitorange!15},
  action1/.style={state, fill=bitblue!15},
  trans/.style={-{Stealth[length=4pt]}, thick},
  >=Stealth
]
  % States for N=4 (8 states)
  \foreach \i in {1,...,4} {
    \node[action0] (s\i) at (\i*1.6, 0) {\i};
  }
  \foreach \i in {5,...,8} {
    \node[action1] (s\i) at (\i*1.6, 0) {\i};
  }

  % Decision boundary
  \draw[dashed, thick, bitgray] (7.2, -1.2) -- (7.2, 1.2)
    node[above, font=\scriptsize\sffamily] {boundary};

  % Action labels
  \node[font=\small\sffamily\bfseries, bitorange] at (4, -1.5)
    {Action 0 (exclude)};
  \node[font=\small\sffamily\bfseries, bitblue] at (10.4, -1.5)
    {Action 1 (include)};

  % Reward arrows (above)
  \draw[trans, bitgreen] (s2.north) to[bend left=30]
    node[above, font=\tiny\sffamily] {reward} (s1.north);
  \draw[trans, bitgreen] (s7.north) to[bend left=30]
    node[above, font=\tiny\sffamily] {reward} (s8.north);

  % Penalty arrows (below)
  \draw[trans, red!70] (s2.south) to[bend right=30]
    node[below, font=\tiny\sffamily] {penalty} (s3.south);
  \draw[trans, red!70] (s7.south) to[bend right=30]
    node[below, font=\tiny\sffamily] {penalty} (s6.south);

  % Boundary crossing
  \draw[trans, red!70] (s5.south) to[bend right=30]
    node[below, font=\tiny\sffamily] {penalty} (s4.south);
\end{tikzpicture}
\caption{A Tsetlin Automaton with $2N = 8$ states. States 1--4 map to
  action~0 (exclude); states 5--8 map to action~1 (include). Rewards
  push toward the ends; penalties push toward the centre boundary.
  A penalty at state~5 crosses the boundary to state~4, flipping the
  action from 1 to~0.}
\label{fig:automaton}
\end{figure}


% ----------------------------------------------------------------------------
\section{Why the Automaton Works: Noise Absorption}
\label{sec:noise-absorption}
% ----------------------------------------------------------------------------

The key property of the Tsetlin Automaton is its ability to resist
noise. Consider an automaton in state~1 (deep in the action~0 region).
Even if it receives a few spurious penalty signals, it must be pushed
through states 2, 3, \ldots, $N$, $N{+}1$ before it changes its
action. The parameter $N$ controls this \concept{noise margin}: larger
$N$ means the automaton needs more consecutive penalties to change its
mind.

\begin{keyinsight}
The number of states $N$ per action controls the trade-off between
\emph{stability} (resistance to noise) and \emph{agility} (speed of
adaptation). $N = 1$ gives a memoryless flip-flop that changes on
every penalty. Large $N$ gives a conservative automaton that changes
slowly but reliably.
\end{keyinsight}

\begin{example}
With $N = 100$ and the automaton in state~1, it takes at least 100
consecutive penalties with no intervening rewards to flip the
action. If the correct action is truly~0, the environment will
produce more rewards than penalties, and the state will drift toward
1 (maximum confidence). If the correct action is truly~1, the
accumulated penalties will eventually push the state past 100 to the
action~1 side, where rewards will then anchor it.
\end{example}


% ----------------------------------------------------------------------------
\section{Convergence Guarantees}
\label{sec:convergence}
% ----------------------------------------------------------------------------

Tsetlin~\cite{tsetlin1961} proved the following result (stated informally):

\begin{theorem}[Tsetlin convergence, informal]
\label{thm:tsetlin-convergence}
If one action produces rewards more frequently than the other, a
two-action Tsetlin Automaton with sufficiently many states will
converge to the better action with probability approaching~1 as
$N \to \infty$.
\end{theorem}

More precisely, let $r_0$ be the reward probability for action~0 and
$r_1$ the reward probability for action~1. If $r_1 > r_0$, then the
probability that the automaton is in an action~1 state converges to~1
as $N \to \infty$.

\begin{engineernote}
In practice, $N$ between 50 and 200 works well for MNIST. The
convergence is fast enough that the automaton settles within a few
training epochs (passes through the data). The default in most
implementations is $N = 100$, giving 200 states per automaton.
\end{engineernote}

The convergence does \emph{not} require independence between
successive feedback signals~\cite{zhang2021convergence}. This is important because in a Tsetlin
Machine, the feedback to each automaton depends on the current input
pattern and the states of other automata---it is highly correlated.
Tsetlin's original proof accommodates this, which is one reason the
framework is robust.


% ----------------------------------------------------------------------------
\section{Multiple Automata Working Together}
\label{sec:team}
% ----------------------------------------------------------------------------

A single automaton makes a single binary decision: include or exclude.
The power of the Tsetlin Machine~\cite{granmo2018tsetlin} comes from assembling many automata
into a \concept{team} that collectively builds a pattern detector (a
clause).

Consider a clause with 10 input features ($x_1, \ldots, x_5$ and
their negations $\bar{x}_1, \ldots, \bar{x}_5$). We assign one
automaton to each of the 10 literals. Each automaton independently
decides whether its literal is ``included'' (action~1) or
``excluded'' (action~0) from the clause.

\begin{figure}[htbp]
\centering
\begin{tikzpicture}[
  auto/.style={draw, rounded corners, minimum width=1.4cm,
    minimum height=0.7cm, font=\scriptsize\sffamily, inner sep=2pt},
  inc/.style={auto, fill=bitblue!15},
  exc/.style={auto, fill=bitgray!10},
  >=Stealth
]
  % Literals
  \foreach \i/\lit/\act in {
    0/$x_1$/include,
    1/$\bar{x}_1$/exclude,
    2/$x_2$/include,
    3/$\bar{x}_2$/include,
    4/$x_3$/exclude
  } {
    \pgfmathsetmacro{\xpos}{\i * 2.5}
    \ifnum\i=1
      \node[exc] (a\i) at (\xpos, 0) {\lit: \act};
    \else\ifnum\i=4
      \node[exc] (a\i) at (\xpos, 0) {\lit: \act};
    \else
      \node[inc] (a\i) at (\xpos, 0) {\lit: \act};
    \fi\fi
  }

  % Clause output
  \node[draw, rounded corners, fill=bitgreen!15, minimum width=2cm,
    font=\small\sffamily\bfseries] (clause) at (5, -1.5)
    {$C = x_1 \wedge \bar{x}_2$};

  % Note
  \node[font=\scriptsize\sffamily, bitgray] at (5, -2.3)
    {(only included literals appear in the clause)};

  % Arrows from included automata to clause
  \draw[->, thick] (a0.south) -- (clause);
  \draw[->, thick] (a2.south) -- (clause);
  \draw[->, thick] (a3.south) -- (clause);
\end{tikzpicture}
\caption{A team of automata forming a clause. Each automaton controls
  one literal. The clause is the AND of all included literals.
  Here, automata for $x_1$, $x_2$, and $\bar{x}_2$ chose ``include,''
  so $C = x_1 \wedge x_2 \wedge \bar{x}_2$. (In practice, including
  both $x_2$ and $\bar{x}_2$ would make the clause always false---the
  training algorithm learns to avoid this.)}
\label{fig:automata-team}
\end{figure}

The clause output is then:
\begin{equation}
  C = \bigwedge_{\substack{k : \text{automaton}_k \\ \text{says include}}} \ell_k
  \label{eq:clause-automata}
\end{equation}

The training algorithm (\cref{ch:clauses}) will adjust the automata
states so that useful literals are included and irrelevant ones are
excluded.


% ----------------------------------------------------------------------------
\section{The Automaton as a Memory Device}
\label{sec:memory}
% ----------------------------------------------------------------------------

It is useful to think of the automaton's state as a form of memory
that accumulates evidence. Each reward strengthens the current
decision; each penalty weakens it. The ``depth'' of memory (controlled
by $N$) determines how much evidence is needed before the automaton
changes its mind.

This is fundamentally different from gradient-based learning:
\begin{itemize}
  \item There is no gradient computation---just integer
    increment/decrement.
  \item There is no learning rate to tune---the step size is always
    $\pm 1$ state.
  \item There is no loss landscape to navigate---the automaton simply
    drifts in the direction that gets more rewards.
\end{itemize}

\begin{keyinsight}
The Tsetlin Automaton replaces all of gradient descent's
machinery---differentiation, learning rates, optimisers---with a
single operation: increment or decrement an integer. This extreme
simplicity is what makes the Tsetlin Machine suitable for
hardware implementation and embedded deployment~\cite{wheeldon2020pervasive}.
\end{keyinsight}


% ----------------------------------------------------------------------------
\section{Stochastic Feedback and Exploration}
\label{sec:stochastic}
% ----------------------------------------------------------------------------

In the Tsetlin Machine, feedback is often made
\concept{stochastic}---that is, the reward or penalty is applied with
some probability rather than deterministically. This serves two
purposes:

\begin{enumerate}
  \item \textbf{Exploration:} Even when a clause is performing well,
    random penalties occasionally push automata toward the boundary,
    allowing the clause to ``try out'' different literal combinations.
  \item \textbf{Diversity:} If all clauses received identical
    deterministic feedback, they would converge to the same pattern.
    Stochastic feedback introduces variation, ensuring that different
    clauses discover different patterns.
\end{enumerate}

The probability of feedback is controlled by a parameter $s$, called
the \concept{specificity} parameter. We will define it precisely in
\cref{ch:clauses}, but the intuition is:

\begin{itemize}
  \item Large $s$ $\Rightarrow$ high probability of reinforcing
    included literals $\Rightarrow$ clauses become more specific
    (include more literals, match fewer inputs).
  \item Small $s$ $\Rightarrow$ lower reinforcement probability
    $\Rightarrow$ clauses stay general (fewer literals, match more
    inputs).
\end{itemize}

\begin{engineernote}
The specificity parameter $s$ is the single most important
hyperparameter of the Tsetlin Machine. Typical values for MNIST are
$s = 3.0$ to $10.0$. A good starting point is $s = 5.0$.
\end{engineernote}


% ----------------------------------------------------------------------------
\section{Summary}
\label{sec:ch5-summary}
% ----------------------------------------------------------------------------

The Tsetlin Automaton is a finite-state machine with a simple update
rule: rewards push the state toward confidence, penalties push it
toward uncertainty. Its key properties are:

\begin{enumerate}
  \item \textbf{Noise tolerance:} The $N$-state memory buffer filters
    out random fluctuations.
  \item \textbf{Guaranteed convergence:} With enough states, the
    automaton provably finds the better action.
  \item \textbf{No calculus:} Learning requires only integer
    increment and decrement.
  \item \textbf{Composability:} Teams of automata can collectively
    learn complex pattern detectors (clauses).
\end{enumerate}

In the next chapter, we will see how the Tsetlin Machine orchestrates
hundreds of these automata, organised into clauses, and trains them
using a clever feedback scheme called Type~I and Type~II feedback.
